% ECONOMICS_PROBLEM: {"id": "G28-66", "title": "Credible bank resolution", "jel_code": "G28", "source_id": "OP-0327"}
% !TeX program = xelatex
\documentclass[11pt,a4paper]{article}
\usepackage[margin=23mm]{geometry}
\usepackage{fontspec}
\setmainfont{Latin Modern Roman}
\usepackage{amsmath,amssymb,mathtools}
\usepackage{xcolor,enumitem,booktabs,longtable,array,xurl,hyperref}
\definecolor{muted}{HTML}{53636C}
\hypersetup{colorlinks=true,urlcolor=blue,linkcolor=blue}
\setlength{\parindent}{0pt}
\setlength{\parskip}{5pt}
\newcommand{\R}{\mathbb R}
\newcommand{\E}{\mathbb E}
\newcommand{\N}{\mathbb N}
\newcommand{\ind}{\mathbf 1}
\newcommand{\Var}{\operatorname{Var}}
\newcommand{\Cov}{\operatorname{Cov}}
\newcommand{\supp}{\operatorname{supp}}
\DeclareMathOperator*{\argmin}{arg\,min}
\DeclareMathOperator*{\argmax}{arg\,max}
\newcommand{\entrymeta}[3]{{\sffamily\small\color{muted}\textbf{\texttt{#1}}\quad|\quad #2\par #3\par}}
\newcommand{\Problemref}[1]{\texttt{#1}}
\newcommand{\JELref}[2]{\texttt{#2} (JEL \texttt{#1})}
\begin{document}
\section*{Credible bank resolution}
\textbf{Problem ID:} \texttt{G28-66}\quad \textbf{JEL:} \texttt{G28}\par
% BEGIN ECONOMICS STATEMENT
\label{problem:OP-0327}
{\sffamily\small\color{muted}Original subject group: Banking, credit and financial stability\par}
\label{definitions:problem:30007609}
\textbf{Audit classification:} Published research agenda. Both sources are real; the empirical 2026 paper does not state this exact welfare-optimal credibility problem. Tucker supplies a broader policy agenda. Missing author/year data are repaired; the sequential policy model is editorial and strategic.
\subsection*{Definitions and precise research target}
Fix a banking system with insured deposits, unsecured claims of stated seniorities, equity and cross-bank exposures. A resolution rule \(a\) specifies loss allocation, recapitalization, transfer of critical functions and temporary funding. Investors price claims before stress anticipating both \(a\) and authorities' possible deviation. A rule is credible if executing it is sequentially optimal for the authority at each specified failure state, given continuation losses and its objective. Let \(W(a)\) include household surplus, functioning payments and credit, fiscal resource costs and spillovers. The task is to find feasible rules maximizing ex ante welfare subject to credibility:
\[A_{\rm cred}=\{a\in A:a\text{ satisfies sequential policy optimality}\},\qquad W_*=\sup_{a\in A_{\rm cred}}W(a).\]
Define \(A\) using insolvency law, available loss-absorbing claims and operational constraints. Model creditor runs and other banks' losses from write-downs rather than assuming bail-in is harmless. Identify credibility from claim pricing or behavior only under explicit assumptions about common risk premia and expected losses. Compare idiosyncratic and systemic failure states. Report whether a rule remains credible when multiple banks fail together. A formal legal power to impose losses and a market expectation that it will be used are different objects.

\textbf{Additional formal definitions and scope (editorial).} Let \(s\) be a failure state, \(D(s)\) feasible authority deviations, and \(V^{A}(d;s,a)\) its conditional continuation value, including fiscal and systemic costs and anticipated future behavior. Credibility of rule \(a\) means \(V^{A}(a(s);s,a)\geq V^{A}(d;s,a)\) for every feasible \(d\in D(s)\) and every state in the declared support, not merely on a historical path. The ex ante feasible set includes this inequality and the legal loss waterfall. Bail-in means contractual conversion or write-down of designated creditor claims; distinguish AT1 contractual loss absorption, resolution bail-in debt and senior liabilities. Price spreads identify risk-neutral expected losses jointly with recovery, stress probability and discounting assumptions; they do not directly observe sequential authority optimality. The time-consistency formulation is an editorial strategic policy model, not a theorem quoted from the empirical bail-in paper. For a nonempty feasible set and a finite optimal value, the design target is an \(\varepsilon\)-optimal feasible rule for each specified \(\varepsilon>0\): its cost is at most the infimum plus \(\varepsilon\), or its welfare is at least the supremum minus \(\varepsilon\). Report infeasibility or an unbounded value separately. An exact optimizer is requested only under additional attainment assumptions; it need not exist for the general rule class.
\subsection*{Variants and assumption changes}
\begin{enumerate}
\item \textbf{Commitment to resolution (editorial).} Assume the loss waterfall can be irrevocably enforced; optimize ex ante welfare without a discretionary credibility constraint.
\item \textbf{Sequential credibility (editorial; strategic policy).} Permit deviations at idiosyncratic and joint bank-failure states; characterize credible rules and identify market-implied expected loss allocation separately.
\end{enumerate}
\subsection*{References and source audit}
\begin{enumerate}
\item Official BIS WP 1356 and PDF verify Alessandro Di Stefano, Yvan Lengwiler and Kumar Rishabh, June 5, 2026; no DOI assigned. Locator: Abstract and Section 7, printed p.40 (PDF p.42). Support: Empirical credibility evidence; no optimal-rule open theorem. Full PDF accessible. URL: \url{https://www.bis.org/publications/working-paper-1356-credibility-bail-in}.
\item Official BIS WP 792 verifies Paul Tucker, 2019; cover June 2019, release July 4, 2019. Comments are separate contributions. Locator: Section 5, printed pp.12--16; Section 7, printed pp.23--24. Support: Financial-resilience policy agenda. Full paper and separate comments accessible. URL: \url{https://www.bis.org/publications/working-paper-792-financial-system-sufficiently-resilient-research-programme-and-policy-agenda}.
\end{enumerate}
\begin{enumerate}\raggedright
\item\hypertarget{definitions:source:30007609:1}{} Alessandro Di Stefano; Yvan Lengwiler; Kumar Rishabh. 2026. \emph{The credibility of bail-in}. Abstract and Section 7, printed p.40 (PDF p.42).
\url{https://www.bis.org/publications/working-paper-1356-credibility-bail-in}.
\item\hypertarget{definitions:source:30007609:2}{} Paul Tucker. 2019. \emph{Is the financial system sufficiently resilient: a research programme and policy agenda}. Section 5, printed pp.12--16; Section 7, printed pp.23--24.
\url{https://www.bis.org/publications/working-paper-792-financial-system-sufficiently-resilient-research-programme-and-policy-agenda}.
\end{enumerate}

\subsection*{Common conventions: Source mathematical conventions}

These conventions apply to each entry unless explicitly replaced.

\textbf{Models and identification.} A primitive model \(m\) specifies all probability laws, feasible choices, information, equilibrium and measurement rules used by its target. The admissible class \(\mathcal M\) is fixed independently of outcomes. If \(P_O(m)\) is its observable probability law and \(\tau(m)\) the target, its identified set at observable law \(P\) is
\[
 \mathcal I_\tau(P)=\{\tau(m):m\in\mathcal M,\ P_O(m)=P\}.
\]
Point identification means that this set is a singleton. An empty set signals incompatibility of data and assumptions. Coordinate bounds need not describe the joint set. Bounds are sharp only if every retained target is attainable or arbitrarily approachable by compatible models. Finite bounds require explicit restrictions on unobserved quantities; unrestricted confounding, hidden wealth, extrapolation or arbitrary equilibrium selection can yield uninformative sets.

\textbf{Causal interventions.} A potential outcome \(Y_i(p)\) is the outcome under a complete policy regime \(p\), including timing, financing, information and market spillovers. The target population and units are fixed before treatment. With interference, use the complete assignment vector or a justified exposure mapping; individual no-interference notation alone cannot identify an economy-wide intervention. A difference of mean potential outcomes does not require the cross-world joint distribution. Conditioning on post-treatment migration, survival, enrollment or employment changes the estimand and needs separate justification.

\textbf{Statistical statements.} An estimator, confidence set, forecast or test specifies a sampling law, model class and loss/coverage criterion. Uniform \((1-\alpha)\)-coverage means \(\inf_{m\in\mathcal M}\Pr_m\{\tau(m)\in C_n\}\geq1-\alpha\), or an explicitly asymptotic version. The whole parameter space trivially has coverage, so informativeness requires a stated width, risk or power objective. A forecasting improvement is relative to a named benchmark, information set and evaluation population.

\textbf{Equilibrium and optimization.} An equilibrium comprises feasible plans, optimality, beliefs where needed, budget constraints and all market/resource clearing. The primitives or a stated benchmark define each correspondence \(\mathcal E(p;m)\). Nonemptiness is assumed or proved, never inferred just from compact policy sets. A selected-equilibrium target uses a declared selection rule; a robust target quantifies over all permitted equilibria. Use suprema or infima unless attainment is established. An \(\varepsilon\)-optimal policy is feasible and within \(\varepsilon\) of the finite optimum value. Report an empty feasible set separately.

\textbf{Welfare and accounting.} Normative utility, interpersonal normalization, social weights, numeraire, discounting and the use of public receipts are primitives. Behavioral utility need not coincide with the welfare criterion. Household welfare includes ownership income and rents once; adding profits again to compensating variation generally double-counts them. A monetary equivalent is defined by an explicit transfer and price convention, with monotonicity and range conditions. Average effects, mechanism contributions, transfers and welfare losses are different targets. Infinite sums require convergence and dynamic feasible plans require terminal or no-Ponzi conditions.

\textbf{Source versions.} A working-paper date, revision date and journal publication year are distinguished. A DOI for a journal version is not automatically the DOI of an earlier manuscript. Locators refer to the stated version and printed pages where specified. A metadata-only check does not verify a claim about a full-text conclusion. Every entry's source subsection narrows the provenance claim accordingly.
% END ECONOMICS STATEMENT
\end{document}
